Un objet de masse $\qty{500}{\g}$ est lancé, vers le haut, d'un point $A$
suivant la ligne de plus grande pente d'un plan incliné faisant un angle
$\alpha=\ang{30}$ avec l'horizontale. Il monte jusqu'en un point $B$, s'arrête
puis repart jusqu'en $A$.
\begin{enumerate}
    \item Faire un schéma.
    \item Les frottements seront d'abord négligés dans un premier temps.

          Faire le bilan des forces qui s'exercent sur l'objet et les
          représenter.
    
          Calculer les travaux effectués par ces forces lors de la montée $A B$,
          lors de la descente $B A$ et lors de l'aller-retour.
    \item On considère maintenant que les frottements ne sont plus négligeables.
          On les assimilera à une force constante de valeur $f=\qty{1.5}{\N}$.

          Faire le bilan des forces qui s'exercent sur l'objet, lors de la
          montée puis lors de la descente, et les représenter.
    
          Calculer les travaux effectués par ces forces lors de la montée $A B$,
          lors de la descente $B A$ et lors de l'aller-retour.
\end{enumerate}

\begin{donnees}
    \begin{itemize}
        \item $A B=\qty{0.80}{\m}$.
        \item $g=\qty{9.8}{\N\per\kg}$.
    \end{itemize}
\end{donnees}

